% ##############################################################################
\section{NoesisMarket: the auction mechanism}
\label{sec:noesis_market}

The auction described throughout this section is a specific
instantiation of the pluggable matching-engine component of
Section~\ref{sec:modularity}; nothing below depends on the matching engine
being a call auction rather than, e.g., a continuous double auction.

A possible implementation of \NoesisMarket{} runs a periodic call auction that
turns a pool of bids and asks into a set of intelligence contracts
(Definition~\ref{def:contract}). This section defines auction rounds and tier
bucketing, the compatibility relation between a bid and an ask, the clearing
problem, the reputation and pricing feedback loop, and the mechanism-design
risks the design has to withstand. 

% ==============================================================================
\subsection{Auction rounds and tier bucketing}
\label{sec:rounds}

Every $T$ minutes (e.g., $T = 5$ minutes), \NoesisMarket{} closes the open
order book and runs one clearing round per capability tier $c \in \calC$. For
clarity of exposition we describe the case in which a bucket for tier $c$
contains exactly the bids with $c_\beta = c$ and the asks with $c_\alpha = c$;
Remark~\ref{rem:tier_generalization} below discusses the natural generalization
in which higher-tier asks are also eligible to fill lower-tier demand.

\begin{example}[Basic auction round]
\label{ex:basic_round}
A buyer submits a bid for 10,000 tasks at \texttt{frontier} capability,
under 2 seconds of latency, at 99.9\% reliability, with a limit price
$p_\beta$. Several sellers submit asks at the \texttt{frontier} tier, with
various typical latencies, reliabilities, and limit prices $p_\alpha$. At
the close of the round, \NoesisMarket{} finds the compatible asks
(Section~\ref{sec:compatibility}), clears the \texttt{frontier} bucket at a
single uniform price, and matches the buyer's bid against one or more of
the compatible asks, up to 10,000 tasks.
\end{example}

Figure~\ref{fig:basic_round} illustrates this example: the bid and the
compatible asks at the \texttt{frontier} tier all flow into \NoesisMarket's
clearing step for that bucket, which produces a single uniform price $p^*$ and
matches the bid against compatible asks up to its 10,000-task limit.

% rendered_images:begin
% ```graphviz
% digraph BasicRound {
%   bgcolor="transparent";
%   pad="0.10";
%   splines=spline;
%   nodesep=0.28;
%   ranksep=0.55;
%   rankdir=LR;
% 
%   node [shape=box,
%         style="rounded,filled",
%         fillcolor="#FFFFFF",
%         color="#9AA9B8",
%         penwidth=1.6,
%         fontname="Helvetica",
%         fontcolor="#243B53",
%         fontsize=10,
%         margin="0.18,0.10",
%         height=0.42];
% 
%   edge [color="#A3B1C0",
%         penwidth=1.2,
%         arrowhead=vee,
%         arrowsize=0.7,
%         fontname="Helvetica",
%         fontsize=8.5,
%         fontcolor="#7B8794"];
% 
%   Bid [label=<<b>Bid</b><br/><font point-size="8.5" color="#A34F6C">10,000 tasks @ frontier<br/>&#8804;2s, &#8805;99.9%, limit p<sub>&#946;</sub></font>>,
%        fillcolor="#F6E1E8", color="#D98CA8", fontcolor="#6B2A44"];
% 
%   subgraph cluster_asks {
%     label     = "Asks at the frontier tier";
%     labelloc  = "t";
%     fontname  = "Helvetica-Bold";
%     fontsize  = 9.5;
%     fontcolor = "#3A4A5C";
%     style     = "rounded,filled";
%     fillcolor = "#F7F9FB";
%     color     = "#C7D0DA";
%     penwidth  = 1.0;
%     margin    = 12;
% 
%     Ask1 [label=<Seller 1: limit p<SUB>&#945;1</SUB>>, fillcolor="#FBEBD4", color="#D9A85F", fontcolor="#6B4517"];
%     Ask2 [label=<Seller 2: limit p<SUB>&#945;2</SUB>>, fillcolor="#FBEBD4", color="#D9A85F", fontcolor="#6B4517"];
%     Ask3 [label=<Seller 3: limit p<SUB>&#945;3</SUB>>, fillcolor="#FBEBD4", color="#D9A85F", fontcolor="#6B4517"];
%   }
% 
%   Clear [label=<<b>NoesisMarket</b><br/><font point-size="8.5" color="#41719C">Clears frontier bucket<br/>at uniform price p*</font>>,
%          fillcolor="#D3E3F3", color="#7CA6CE", fontcolor="#1F4E79"];
%   Matched [label=<<b>Matched</b><br/><font point-size="8.5" color="#2E7C74">Up to 10,000 tasks @ p*</font>>,
%            fillcolor="#D4EEEB", color="#7FC2B9", fontcolor="#1F5A54"];
% 
%   Bid -> Clear [label="bid"];
%   Ask1 -> Clear;
%   Ask2 -> Clear [label="compatible asks"];
%   Ask3 -> Clear;
%   Clear -> Matched;
% }
% ```
% label=fig:basic_round
% caption=The basic auction round of Example~\ref{ex:basic_round}: a bid and compatible asks at the frontier tier are cleared into a match at a single uniform price.
% rendered_images:end
% render_images:begin
\begin{figure}[H]
  \includegraphics[width=\linewidth]{04_noesis_market.tex.figs/04_noesis_market.1.png}
  \caption{The basic auction round of Example~\ref{ex:basic_round}: a bid and compatible asks at the frontier tier are cleared into a match at a single uniform price.}
  \label{fig:basic_round}
\end{figure}
% render_images:end

\begin{remark}[Tiered commodity pricing]
\label{rem:tiered_pricing}
Because each tier clears independently, prices can diverge across tiers
within the same round, in the same way that peak and off-peak electricity
prices diverge within the same day. A \texttt{cheap} tier and a
\texttt{frontier} tier are not substitutes from the buyer's perspective and
are not expected to clear at the same price.
\end{remark}

Figure~\ref{fig:auction_timeline} lays out this timeline: the order book
accepts bids and asks continuously, and every $T$ minutes the round closes
and a separate uniform-price clearing is computed independently per
capability-tier bucket, so \texttt{cheap} and \texttt{frontier} can clear
at different prices in the same round (Remark~\ref{rem:tiered_pricing}).

% rendered_images:begin
% ```tikz
% \definecolor{orderband}{RGB}{251,235,212}
% \definecolor{timelinecolor}{RGB}{133,146,163}
% \definecolor{dividercolor}{RGB}{199,208,218}
% \definecolor{cheapfill}{RGB}{211,227,243}
% \definecolor{cheapline}{RGB}{124,166,206}
% \definecolor{cheaptext}{RGB}{31,78,121}
% \definecolor{frontfill}{RGB}{233,225,247}
% \definecolor{frontline}{RGB}{168,143,217}
% \definecolor{fronttext}{RGB}{69,41,107}
% \definecolor{bidcolor}{RGB}{163,79,108}
% \definecolor{askcolor}{RGB}{148,106,46}
% \fontfamily{phv}\selectfont
% \fill[orderband, opacity=0.55] (0,-0.3) rectangle (12.8,0.3);
% \node[align=center] at (6.4,0) {\footnotesize order book open: continuously accepting bids / asks};
% \draw[very thick, ->, >=stealth, timelinecolor] (0,-1.0) -- (13.4,-1.0);
% \draw[dashed, dividercolor] (2,-1.15) -- (2,2.6);
% \draw[dashed, dividercolor] (6.4,-1.15) -- (6.4,2.6);
% \draw[dashed, dividercolor] (10.8,-1.15) -- (10.8,2.6);
% \node[below] at (2,-1.15) {\footnotesize round closes at $t$};
% \node[below] at (6.4,-1.15) {\footnotesize round closes at $t+T$};
% \node[below] at (10.8,-1.15) {\footnotesize round closes at $t+2T$};
% % Bids arriving from the top and asks arriving from the bottom, into the
% % order band, at a few illustrative points along the timeline.
% \foreach \x in {0.8, 3.1, 5.4, 7.7, 9.9, 12.1} {
%   \draw[->, >=stealth, thin, bidcolor] (\x,0.75) -- (\x,0.32);
% }
% \node[font=\tiny, text=bidcolor] at (0.8,0.92) {bid};
% \foreach \x in {1.6, 3.9, 6.2, 8.5, 10.7, 12.9} {
%   \draw[->, >=stealth, thin, askcolor] (\x,-0.75) -- (\x,-0.32);
% }
% \node[font=\tiny, text=askcolor] at (1.6,-0.92) {ask};
% \node[draw=cheapline, rounded corners, fill=cheapfill, minimum width=1.9cm, minimum height=0.55cm, font=\footnotesize, align=center, text=cheaptext] at (2,1.1) {cheap tier\\clears at $p^*_{\text{cheap}}$};
% \node[draw=frontline, rounded corners, fill=frontfill, minimum width=1.9cm, minimum height=0.55cm, font=\footnotesize, align=center, text=fronttext] at (2,2.0) {frontier tier\\clears at $p^*_{\text{frontier}}$};
% \node[draw=cheapline, rounded corners, fill=cheapfill, minimum width=1.9cm, minimum height=0.55cm, font=\footnotesize, align=center, text=cheaptext] at (6.4,1.1) {cheap tier\\clears at $p^*_{\text{cheap}}$};
% \node[draw=frontline, rounded corners, fill=frontfill, minimum width=1.9cm, minimum height=0.55cm, font=\footnotesize, align=center, text=fronttext] at (6.4,2.0) {frontier tier\\clears at $p^*_{\text{frontier}}$};
% \node[draw=cheapline, rounded corners, fill=cheapfill, minimum width=1.9cm, minimum height=0.55cm, font=\footnotesize, align=center, text=cheaptext] at (10.8,1.1) {cheap tier\\clears at $p^*_{\text{cheap}}$};
% \node[draw=frontline, rounded corners, fill=frontfill, minimum width=1.9cm, minimum height=0.55cm, font=\footnotesize, align=center, text=fronttext] at (10.8,2.0) {frontier tier\\clears at $p^*_{\text{frontier}}$};
% ```
% label=fig:auction_timeline
% caption=The \NoesisMarket{} auction timeline.
% rendered_images:end
% render_images:begin
\begin{figure}[H]
  \includegraphics[width=\linewidth]{04_noesis_market.tex.figs/04_noesis_market.2.png}
  \caption{The \NoesisMarket{} auction timeline.}
  \label{fig:auction_timeline}
\end{figure}
% render_images:end

% ==============================================================================
\subsection{Compatibility}
\label{sec:compatibility}

Within a tier bucket, a bid is not compatible with every ask: the bid's
latency and reliability requirements still need to be met.

\begin{definition}[Compatibility]
\label{def:compatibility}
A bid $\beta = (n_\beta, c_\beta, \ell_\beta, r_\beta, p_\beta)$ and an ask
$\alpha = (n_\alpha, c_\alpha, \ell_\alpha, r_\alpha, p_\alpha)$ are
\emph{compatible}, written $\beta \sim \alpha$, if
$$
  c_\alpha \succeq c_\beta,
  \qquad
  \ell_\alpha \leq \ell_\beta,
  \qquad
  r_\alpha \geq r_\beta.
$$
\end{definition}

Compatibility is a feasibility filter, not a scoring function: capability,
latency, and reliability are treated as hard constraints rather than as
continuous inputs to the price. This is a deliberate simplification. It
avoids having to define an exchange rate between, say, one unit of latency
and one unit of price, at the cost of potentially leaving compatible
volume unmatched when a bid's requirements are only narrowly missed;
Section~\ref{sec:open_questions} discusses relaxing this into a scored
compatibility measure.

Restricted to a tier bucket $c$, compatibility defines a bipartite graph
$G_c = (B_c \cup A_c, E_c)$ between the bids $B_c$ and asks $A_c$ in the
bucket, with an edge $(\beta, \alpha) \in E_c$ whenever $\beta \sim
\alpha$.

Figure~\ref{fig:compatibility_graph} shows a small example of such a
graph: an edge marks a bid/ask pair whose capability, latency, and
reliability terms are mutually satisfiable, and bid 4 has no compatible
ask in this bucket and is left unmatched regardless of price.

% rendered_images:begin
% ```graphviz
% digraph Compatibility {
%   // ---------------------------------------------------------------
%   bgcolor="transparent";
%   pad="0.06";
%   splines=line;
%   nodesep=0.28;
%   ranksep=0.75;
%   rankdir=TB;
% 
%   node [shape=box,
%         style="rounded,filled",
%         fillcolor="#FFFFFF",
%         color="#9AA9B8",
%         penwidth=1.6,
%         fontname="Helvetica",
%         fontcolor="#243B53",
%         fontsize=9,
%         margin="0.14,0.08",
%         height=0.35];
% 
%   edge [color="#A3B1C0",
%         penwidth=1.1,
%         arrowhead=none,
%         arrowsize=0.6,
%         fontname="Helvetica",
%         fontsize=7,
%         fontcolor="#7B8794"];
% 
%   // ---------------------------------------------------------------
%   // Figure: asks at the top, bids at the bottom
%   // ---------------------------------------------------------------
%   // bid (demand-side) : rose   |   ask (supply-side) : warm sand
% 
%   subgraph cluster_asks {
%     label     = "Asks A_c";
%     labelloc  = "t";
%     fontname  = "Helvetica-Bold";
%     fontsize  = 9.5;
%     fontcolor = "#3A4A5C";
%     style     = "rounded,filled";
%     fillcolor = "#F7F9FB";
%     color     = "#C7D0DA";
%     penwidth  = 1.0;
%     margin    = 10;
%     rank = same;
% 
%     A1 [label="ask 1", fillcolor="#FBEBD4", color="#D9A85F", fontcolor="#6B4517"];
%     A2 [label="ask 2", fillcolor="#FBEBD4", color="#D9A85F", fontcolor="#6B4517"];
%     A3 [label="ask 3", fillcolor="#FBEBD4", color="#D9A85F", fontcolor="#6B4517"];
%   }
%   subgraph cluster_bids {
%     label     = "Bids B_c";
%     labelloc  = "t";
%     fontname  = "Helvetica-Bold";
%     fontsize  = 9.5;
%     fontcolor = "#3A4A5C";
%     style     = "rounded,filled";
%     fillcolor = "#F7F9FB";
%     color     = "#C7D0DA";
%     penwidth  = 1.0;
%     margin    = 10;
%     rank = same;
% 
%     B1 [label="bid 1", fillcolor="#F6E1E8", color="#D98CA8", fontcolor="#6B2A44"];
%     B2 [label="bid 2", fillcolor="#F6E1E8", color="#D98CA8", fontcolor="#6B2A44"];
%     B3 [label="bid 3", fillcolor="#F6E1E8", color="#D98CA8", fontcolor="#6B2A44"];
%     B4 [label="bid 4", fillcolor="#F6E1E8", color="#D98CA8", fontcolor="#6B2A44"];
%   }
% 
%   A1 -> B1;
%   A2 -> B1;
%   A2 -> B2;
%   A3 -> B2;
%   A3 -> B3;
% }
% ```
% label=fig:compatibility_graph
% caption=A small example compatibility graph $G_c$ within one tier bucket, asks at the top and bids at the bottom.
% rendered_images:end
% render_images:begin
\begin{figure}[H]
  \includegraphics[width=\linewidth]{04_noesis_market.tex.figs/04_noesis_market.3.png}
  \caption{A small example compatibility graph $G_c$ within one tier bucket, asks at the top and bids at the bottom.}
  \label{fig:compatibility_graph}
\end{figure}
% render_images:end

% ==============================================================================
\subsection{The clearing problem}
\label{sec:clearing_problem}

\begin{problem}[NoesisMarket clearing]
\label{prob:clearing}
Fix a tier $c \in \calC$ and round $t$, with bids $B_c$, asks $A_c$, and
compatibility graph $G_c = (B_c \cup A_c, E_c)$ as above. For each edge
$(\beta, \alpha) \in E_c$, let $x_{\beta \alpha} \geq 0$ denote the number
of tasks matched between $\beta$ and $\alpha$. The \NoesisMarket{}
clearing problem is to find a uniform clearing price
$p^*(c, t) \in \bbR_+$ and nonnegative quantities $\{x_{\beta \alpha}\}$
that maximize the total number of matched tasks
\begin{equation}
  \label{eq:clearing_objective}
  \sum_{(\beta, \alpha) \in E_c} x_{\beta \alpha}
\end{equation}
subject to
\begin{align}
  \label{eq:buyer_quantity}
  \sum_{\alpha \,:\, (\beta, \alpha) \in E_c} x_{\beta \alpha}
    &\leq n_\beta, & \forall \beta \in B_c, \\
  \label{eq:seller_quantity}
  \sum_{\beta \,:\, (\beta, \alpha) \in E_c} x_{\beta \alpha}
    &\leq n_\alpha, & \forall \alpha \in A_c, \\
  \label{eq:price_limit}
  x_{\beta \alpha} &= 0
    \text{ unless } p_\beta \geq p^*(c, t) \geq p_\alpha,
    & \forall (\beta, \alpha) \in E_c.
\end{align}
\end{problem}

Constraints~\eqref{eq:buyer_quantity} and \eqref{eq:seller_quantity}
prevent over-filling a bid or an ask. Constraint~\eqref{eq:price_limit}
enforces the uniform-price property: at most one price clears the bucket,
and only bid/ask pairs whose limit prices bracket that price may be
matched, exactly as in a standard call auction \cite{mcafee1992}. Each
matched bid $\beta$ gives rise to a contract $\kappa$
(Definition~\ref{def:contract}) with
$\ntasksattr = \sum_\alpha x_{\beta \alpha}$, $\clevelattr = c$,
$\lmaxattr = \ell_\beta$, $\rminattr = r_\beta$, and
$\priceattr = p^*(c, t) \cdot \ntasksattr$, dispatched to \NoesisServer{}
for execution.

For fixed $p^*(c, t)$, Problem~\ref{prob:clearing} is a transportation
problem, a special case of linear programming, and is solvable in
polynomial time by standard methods. Choosing $p^*(c, t)$ itself is, in
general, a search over the (finite) set of distinct limit prices submitted
in the bucket, since the optimal matched volume as a function of a
candidate uniform price is piecewise constant.

\begin{remark}[Generalization across tiers]
\label{rem:tier_generalization}
Definition~\ref{def:compatibility} already allows $c_\alpha \succ c_\beta$:
a seller offering a higher tier than requested is compatible with a
lower-tier bid. A full implementation would therefore let a tier-$c$
bucket draw on asks from tier $c$ and above, at the cost of a larger
compatibility graph per round. We do not expect this to be exploited
routinely by rational sellers in practice, since a higher tier's clearing
price is typically at least as high as a lower tier's (Remark~\ref{rem:tiered_pricing}),
so a seller usually has no incentive to divert capacity away from its own
tier's bucket; the case where it is worth doing so, e.g., a temporarily
oversupplied higher tier, is exactly the kind of cross-tier fungibility
that a naive per-tier bucketing could miss, so the possibility is still
worth flagging.
\end{remark}

\begin{proposition}[Existence of a clearing price]
\label{prop:existence}
Consider a tier bucket $c$ in which the compatibility graph $G_c$ is
complete, i.e., every bid in $B_c$ is compatible with every ask in $A_c$
(for instance, because every ask in the bucket meets the most demanding
latency and reliability requirement among the bids in the bucket). Let
$D(p) = \sum_{\beta \in B_c : p_\beta \geq p} n_\beta$ denote cumulative
demand at price $p$ and $S(p) = \sum_{\alpha \in A_c : p_\alpha \leq p}
n_\alpha$ denote cumulative supply at price $p$. Since $D$ is
non-increasing and $S$ is non-decreasing in $p$, there exists a price
$p^*(c, t)$ at which the volume-maximizing matching of
Problem~\ref{prob:clearing} equals $\min(D(p^*), S(p^*))$, and no price
supports strictly more matched volume.
\end{proposition}

\begin{proof}[Proof sketch]
Under a complete compatibility graph, Problem~\ref{prob:clearing} reduces
to a standard uniform-price double auction: for any candidate price $p$,
the maximum feasible matched volume is $\min(D(p), S(p))$, achieved by
matching the highest-limit-price compatible bids against the
lowest-limit-price compatible asks. As $p$ increases, $D(p)$ is
non-increasing and $S(p)$ is non-decreasing, so $\min(D(p), S(p))$ is
maximized at any $p^*$ where the two curves cross (or, in the discrete
case, at the price step immediately before $D$ falls below $S$). This is
the classical Walrasian argument for existence of a market-clearing price
\cite{mcafee1992}.
\end{proof}

Figure~\ref{fig:supply_demand_curves} illustrates
Proposition~\ref{prop:existence}: the non-increasing cumulative demand
curve $D(p)$ and the non-decreasing cumulative supply curve $S(p)$ cross
at the clearing price $p^*(c, t)$, where the matched volume equals
$\min(D(p^*), S(p^*))$.

% rendered_images:begin
% ```tikz
% \definecolor{axiscolor}{RGB}{74,85,104}
% \definecolor{demandcolor}{RGB}{163,79,108}
% \definecolor{supplycolor}{RGB}{65,113,156}
% \definecolor{crosscolor}{RGB}{192,69,91}
% 
% \draw[->, >=stealth, thick, axiscolor] (0,0) -- (8.6,0) node[right, font=\small] {price $p$};
% \draw[->, >=stealth, thick, axiscolor] (0,0) -- (0,6.3) node[above, font=\small] {quantity};
% 
% % Cumulative demand curve D(p), non-increasing (step function).
% \draw[thick, demandcolor]
%   (0,5.6) -- (1.4,5.6) -- (1.4,4.8) -- (2.6,4.8) -- (2.6,3.9) -- (3.6,3.9)
%   -- (3.6,3.0) -- (4.2,3.0) -- (4.2,2.2) -- (5.4,2.2) -- (5.4,1.3)
%   -- (6.6,1.3) -- (6.6,0.4) -- (8,0.4);
% \node[demandcolor, font=\small] at (1.6,5.9) {$D(p)$};
% 
% % Cumulative supply curve S(p), non-decreasing (step function).
% \draw[thick, supplycolor]
%   (0,0.5) -- (1.0,0.5) -- (1.0,1.1) -- (2.0,1.1) -- (2.0,1.8) -- (2.9,1.8)
%   -- (2.9,2.4) -- (3.6,2.4) -- (3.6,3.0) -- (4.2,3.0) -- (4.2,3.7)
%   -- (5.2,3.7) -- (5.2,4.5) -- (6.3,4.5) -- (6.3,5.3) -- (7.5,5.3);
% \node[supplycolor, font=\small] at (6.9,5.55) {$S(p)$};
% 
% % Clearing price p*(c,t) and matched volume.
% \draw[dashed, crosscolor] (4.2,0) -- (4.2,3.0) -- (0,3.0);
% \filldraw[crosscolor] (4.2,3.0) circle (2.2pt);
% \node[below, crosscolor, font=\small] at (4.2,-0.05) {$p^*(c,t)$};
% \node[left, crosscolor, font=\small] at (-0.05,3.0) {$\min(D(p^*),S(p^*))$};
% ```
% label=fig:supply_demand_curves
% caption=The cumulative demand curve $D(p)$ and supply curve $S(p)$ crossing at the clearing price $p^*(c, t)$ (Proposition~\ref{prop:existence}).
% rendered_images:end
% render_images:begin
\begin{figure}[H]
  \includegraphics[width=\linewidth]{04_noesis_market.tex.figs/04_noesis_market.4.png}
  \caption{The cumulative demand curve $D(p)$ and supply curve $S(p)$ crossing at the clearing price $p^*(c, t)$ (Proposition~\ref{prop:existence}).}
  \label{fig:supply_demand_curves}
\end{figure}
% render_images:end

When the compatibility graph is not complete, Problem~\ref{prob:clearing}
remains a well-posed linear program for any fixed candidate price, but a
single uniform price is no longer guaranteed to support a volume-maximizing
matching with zero unmatched compatible pairs; in that case, $p^*(c, t)$ is
a design choice among the prices that support the (still well-defined)
volume-maximizing matching, for instance the limit price of the marginal
cleared ask.

% ==============================================================================
\subsection{Reputation and pricing feedback loop}
\label{sec:reputation}

Every contract $\kappa$ dispatched to \NoesisServer{} for execution is later
reported back with a measured reliability $\hat r(\kappa, W)$ over its
fulfillment window $W$ (Section~\ref{sec:fulfillment_monitoring}).
\NoesisMarket{} maintains a reputation score per seller and uses it to gate
participation in future rounds. As with the matching engine, the specific
update rule given below is one instantiation of the pluggable
reputation-and-feedback component of Section~\ref{sec:modularity}.

\begin{definition}[Seller reputation]
\label{def:reputation}
For a seller whose ask $\alpha$ was matched into contract $\kappa$ at
round $t$, the reputation score $\rho_\alpha(t) \in [0, 1]$ is updated as
\begin{equation}
  \label{eq:reputation_update}
  \rho_\alpha(t + 1)
    = (1 - \lambda) \, \rho_\alpha(t)
    + \lambda \cdot \mathbb{1}\!\left[
        \hat r(\kappa, W) \geq \rminattr(\kappa)
      \right],
\end{equation}
where $\lambda \in (0, 1]$ is a decay rate controlling how quickly recent
performance dominates the score.
\end{definition}

A seller with $\rho_\alpha(t) < \rho_{\min}$, for a market-wide threshold
$\rho_{\min}$, is excluded from submitting asks in the tier where the
shortfall occurred, or is required to submit asks at a lower tier
reflecting the downgraded assessment of its actual capability. This
mirrors the non-performance penalties used in electricity capacity markets
to keep generators accountable for the capacity they committed to sell
\cite{cramtonstoft2005}.

\begin{example}[Default]
\label{ex:default}
A seller wins a contract at the \texttt{frontier} tier with
$\rminattr = 0.999$, but \NoesisServer{} measures $\hat r(\kappa, W) =
0.97$ over the fulfillment window. The violation is reported to
\NoesisMarket, which lowers $\rho_\alpha(t)$ via
Equation~\eqref{eq:reputation_update}; if $\rho_\alpha$ falls below
$\rho_{\min}$, the seller is excluded from the \texttt{frontier} bucket in
subsequent rounds until its reputation recovers.
\end{example}

Figure~\ref{fig:reputation_loop} traces this loop from a cleared contract
back to future eligibility: a contract $\kappa$ dispatched to
\NoesisServer{} yields a measured reliability $\hat r(\kappa, W)$, which
drives the reputation update of Equation~\eqref{eq:reputation_update}, and
the resulting score gates the seller's eligibility to submit asks in
future \NoesisMarket{} rounds, closing the loop.

% rendered_images:begin
% ```graphviz
% digraph ReputationLoop {
%   // ---------------------------------------------------------------
%   bgcolor="transparent";
%   pad="0.10";
%   splines=spline;
%   nodesep=0.45;
%   ranksep=0.45;
%   rankdir=TB;
% 
%   node [shape=box,
%         style="rounded,filled",
%         fillcolor="#FFFFFF",
%         color="#9AA9B8",
%         penwidth=1.8,
%         fontname="Helvetica",
%         fontcolor="#243B53",
%         fontsize=10,
%         margin="0.20,0.12",
%         height=0.50,
%         width=2.6];
% 
%   edge [color="#A3B1C0",
%         penwidth=1.4,
%         arrowhead=vee,
%         arrowsize=0.8,
%         fontname="Helvetica",
%         fontsize=8.5,
%         fontcolor="#7B8794"];
% 
%   // ---------------------------------------------------------------
%   // Figure: five numbered steps, read top to bottom, closed by a
%   // dashed feedback edge back to step 1.
%   // ---------------------------------------------------------------
%   // process : blue   |   artifact : sage green   |   risk / consequence : red
% 
%   Market      [label=<<b>1. NoesisMarket round t</b><br/><font point-size="8.5" color="#41719C">Clears ask α into contract κ</font>>,
%                fillcolor="#D3E3F3", color="#7CA6CE", fontcolor="#1F4E79"];
%   Contract    [label=<<b>2. Matched contract κ</b><br/><font point-size="8.5" color="#4F7A5A">Dispatched to NoesisServer</font>>,
%                fillcolor="#DFEDE0", color="#8FB79A", fontcolor="#2E5A3D"];
%   Measured    [label=<<b>3. Measured reliability</b><br/><font point-size="8.5" color="#4F7A5A">r(κ, W) over window W</font>>,
%                fillcolor="#DFEDE0", color="#8FB79A", fontcolor="#2E5A3D"];
%   Reputation  [label=<<b>4. Reputation update</b><br/><font point-size="8.5" color="#41719C">ρ_α(t+1), Eq. reputation_update</font>>,
%                fillcolor="#D3E3F3", color="#7CA6CE", fontcolor="#1F4E79"];
%   Gate        [label=<<b>5. Gated eligibility</b><br/><font point-size="8.5" color="#A83D3D">Excluded if ρ_α &lt; ρ_min</font>>,
%                fillcolor="#F8DCDC", color="#D97A7A", fontcolor="#6B1F1F"];
% 
%   Market -> Contract [label="clears"];
%   Contract -> Measured [label="fulfillment window W"];
%   Measured -> Reputation [label="compare to R_min(κ)"];
%   Reputation -> Gate [label="threshold ρ_min"];
%   Gate -> Market [
%     label="future asks",
%     style=dashed, penwidth=1.6,
%     color="#C0455B", fontcolor="#C0455B",
%     constraint=false
%   ];
% }
% ```
% label=fig:reputation_loop
% caption=The \NoesisMarket{} reputation and pricing feedback loop, as a five-step flow chart closed by a feedback edge.
% rendered_images:end
% render_images:begin
\begin{figure}[H]
  \includegraphics[width=\linewidth]{04_noesis_market.tex.figs/04_noesis_market.5.png}
  \caption{The \NoesisMarket{} reputation and pricing feedback loop, as a five-step flow chart closed by a feedback edge.}
  \label{fig:reputation_loop}
\end{figure}
% render_images:end

% ==============================================================================
\subsection{Mechanism-design risks}
\label{sec:mechanism_design_risks}

The design above leaves at least three classical mechanism-design problems
open, which we flag here and revisit in
Section~\ref{sec:open_questions}.

\begin{itemize}
  \item \textbf{Capability misrepresentation}: an ask may overstate
    $c_\alpha$ to be eligible for a higher-clearing-price bucket than the
    seller can actually serve. \Noesis{} does not attempt to prevent this
    ex ante; it is deterred ex post through the reputation and pricing
    feedback loop of Section~\ref{sec:reputation}, once fulfillment
    monitoring detects the shortfall.

  \item \textbf{Bid shading}: uniform-price double auctions are not, in
    general, dominant-strategy incentive compatible, so a rational buyer
    or seller may misreport its true valuation \cite{myerson1981}.
    Dominant-strategy mechanisms exist for restricted settings, such as
    single-unit double auctions \cite{mcafee1992}, but do not obviously
    extend to the multi-unit, multi-attribute setting of
    Problem~\ref{prob:clearing} without efficiency loss.

  \item \textbf{Collusion}: a thin tier bucket, with few sellers, is
    vulnerable to the kind of coordinated bidding studied in the auction
    literature on bidding rings \cite{mcafeemcmillan1992}, which is more
    of a concern at the \texttt{frontier} tier, where the number of
    credible providers is small by construction.
\end{itemize}

We do not claim to resolve these problems in this paper; the auction
design in Problem~\ref{prob:clearing} is a starting point whose robustness
to strategic behavior is left for empirical and theoretical follow-up work.
